\section{Who Owns What}
\label{sec:ch06-who-owns-what}
\index{federation!structural directives|(}%

The previous section was about fields that live on two sides of a boundary at
once. This section is about directives that change who is on which side.

\subsection*{\code{@override}: moving a field from one subgraph to another}

\index{federation!override@\code{@override}}%
A field that lives in one subgraph can be migrated to another, and
\code{@override} is how you tell the router it has moved:

\begin{graphqlsdl}
# The first version: Catalog owns `price`
# Catalog subgraph:
type Product @key(fields: "id") {
  id: ID!
  price: Money!
}

# Pricing subgraph:
type Product @key(fields: "id") {
  id: ID!
}

# Later: Pricing takes over `price`
# Pricing subgraph, after the change:
type Product @key(fields: "id") {
  id: ID!
  price: Money! @override(from: "catalog")
}
\end{graphqlsdl}

\index{federation!override progressive migration@\code{@override} progressive migration}%
The \code{from} argument names the subgraph that used to resolve the field. The
router sees it and changes where it sends queries for \code{price}, immediately,
for every client. There is no deprecation window, no dual-resolution period, and
no rollback mechanism beyond deploying the old schema again. That is a sharp
tool, and chapter~\ref{ch:strangling-the-monolith} is about using it safely:
progressive migration, percentage-based traffic shifting, and the moment it is
worth pulling rather than the easier path of keeping the field in two places.

The migrated field has to match in name and type. Pricing cannot decide
\code{price} is now \code{Float} when Catalog declared it \code{Money}. That is
not an override; it is a different field, and composition rejects it.

\subsection*{\code{@interfaceObject}: contributing to an interface from a distance}

\index{federation!interfaceObject@\code{@interfaceObject}}%
\index{entities!entity interfaces}%
An interface can be an entity. Put \code{@key} on an interface, and every object
type implementing it is an entity too, reachable through the interface's key.
Then a subgraph that does not define the interface itself can contribute fields
to every entity implementing it, all at once, without naming each type.

\begin{graphqlsdl}
# In the Media subgraph
interface Media @key(fields: "id") {
  id: ID!
  url: String!
}

type Image implements Media @key(fields: "id") {
  id: ID!
  url: String!
  width: Int!
  height: Int!
}

# In the Search subgraph
type Image @key(fields: "id") {
  id: ID!
}

type Media @interfaceObject @key(fields: "id") {
  id: ID!
  relevanceScore: Float!
}
\end{graphqlsdl}

\index{federation!interfaceObject mechanics@\code{@interfaceObject} mechanics}%
The Search subgraph never declared \code{relevanceScore}
on \code{Image}. It declared it on \code{Media}, marked
as \code{@interfaceObject}, and the composer adds the field to every type
implementing \code{Media}. If a
seventh type implements \code{Media} next month, Search's contribution applies
to it automatically, without a schema change in Search. The subgraph still needs
a reference resolver for each implementing type, because the router still sends
representations, and the resolver has to know which type it is being asked for.

\index{federation!interfaceObject vs Federation v1@\code{@interfaceObject} vs Federation v1}%
This directive replaced a Federation v1 pattern where every implementing type
had to be listed and extended separately. That pattern did not survive a new
implementation arriving, and \code{@interfaceObject} is the fix.

\subsection*{\code{@inaccessible}: hiding things from the public API}

\index{federation!inaccessible@\code{@inaccessible}}%
Not everything a subgraph defines belongs in the supergraph. A field that exists
only so another subgraph can \code{@require} it, or an internal type that
happens to be in the schema because the framework generated it, can be marked
\code{@inaccessible}:

\begin{graphqlsdl}
type Product @key(fields: "id") {
  id: ID!
  internalCatalogId: String! @inaccessible
  title: String!
}
\end{graphqlsdl}

The client never sees \code{internalCatalogId} in the supergraph schema, and
introspection does not show it. The field still exists in the subgraph, and
other subgraphs can still reference it in \code{@requires}. What changes is
that the public API is cleaned of it.

\index{federation!inaccessible types@\code{@inaccessible} types}%
The directive works on types, enum values and arguments too, not only on fields.
A whole type marked \code{@inaccessible} disappears from the supergraph while
its subgraph goes on using it. The most common use is hiding implementation
details that the gateway does not need to know about, and
chapter~\ref{ch:identity-and-authorization} uses it to keep internal
authorisation fields out of the routeable schema.

\subsection*{\code{@tag}: labelling things for tools}

\index{federation!tag@\code{@tag}}%
\code{@tag(name: "string")} attaches a label to any schema element. It does
nothing to composition, query planning or execution. It exists for tooling: a
registry can filter types by tag, a dashboard can colour nodes by tag, a CI
check can fail pull requests that add a field without the right tag. Mosaic uses
it sparingly, because a tag that is not read by a machine is a comment that
looks like a directive, and comments are cheaper. Chapter~\ref{ch:ci-cd-for-a-federated-graph}
is where a machine starts reading them.

\index{federation!structural directives|)}%
